\section{Evaluation}
\label{sec:evaluation}

\subsection{Paired Phase Experiment}

The event simulator implements the contention reduction in
\eqref{eq:reduction}. A transmission occupies \SI{2}{\milli\second}; Wi-Fi and
NR-U use equal contention windows $[15,63]$, and NR-U access is aligned to a
\SI{250}{\micro\second} grid. Exogenous arrival and contention streams are
paired between Adaptive DB-LBT and fixed TMC. The 32 formal seeds, ranging
from 20029 to 20341, were frozen before the pilot and do not overlap the three
discovery seeds.

Each 32,768-round run alternates two 4096-round phases four times. The light
phase activates four Wi-Fi and four NR-U nodes with Poisson rate 0.025
packets/ms per node. The high phase activates all six nodes of each technology
at rate 0.045. Thus each seed contains seven observed transitions and equal
dwell time in both regimes.

\begin{table}[!t]
\caption{Formal Confirmation Design}
\label{tab:formal-design}
\centering
\footnotesize
\begin{tabularx}{\columnwidth}{@{}p{25mm}Y@{}}
\toprule
Item & Frozen value \\
\midrule
Compared policies & Distributed LinUCB over $q_4/q_{20}$; fixed $q_{20}$ TMC \\
Formal evidence & 32 paired seeds; discovery seeds excluded \\
Light phase & $4+4$, Poisson 0.025 packets/ms/node \\
High phase & $6+6$, Poisson 0.045 packets/ms/node \\
Schedule & Eight 4096-round phases; 32,768 rounds/run \\
Inference & Seed-paired bootstrap, 10,000 resamples \\
\bottomrule
\end{tabularx}
\end{table}

The reporting utility is separate from the local reward in
\eqref{eq:reward}:
\begin{align}
U_D^{\rm eval}&=1-\operatorname{clip}(D_{95}/500\,\mathrm{ms},0,1),\nonumber\\
U_{\rm eval}&=(A_{\rm total}+U_D^{\rm eval}+J)/3-0.25P_{\rm col}.
\label{eq:evaluation-utility}
\end{align}
It combines effective data airtime, P95 access delay, Jain fairness, and event
collision probability. We first recompute all four components within every
contiguous phase, average the eight phase utilities within a seed, and form
paired bootstrap intervals across seeds. This phase-conditional estimand keeps
the active population fixed inside each metric calculation.

\begin{figure*}[!t]
  \centering
  \includegraphics[width=.96\textwidth]{restricted-confirmation.pdf}
  \caption{Independent 32-seed confirmation. Left: Adaptive-minus-TMC utility;
  intervals are paired 95\% bootstrap CIs. ``Phase avg.'' is the primary
  within-phase estimator, while ``Mixed-set diag.'' computes one utility over
  the entire run. Right: online choices are restricted to the two profiles in
  Table~\ref{tab:profiles}.}
  \label{fig:restricted}
\end{figure*}

\subsection{Performance and Decisions}

Figure~\ref{fig:restricted} shows a strong regime split. In the light $4+4$
phase, Adaptive DB-LBT improves utility by $0.015735$ (95\% CI
$[0.015322,0.016145]$), and all 32 seed differences are positive. In the high
$6+6$ phase the effect is $-0.001058$
($[-0.001283,-0.000843]$); only one seed is positive. Averaging equal-duration
phases gives $+0.007338$ ($[0.007093,0.007584]$), positive for every seed.

The component differences explain that result. Relative to TMC, the
phase-averaged collision probability falls by 0.032288 and Jain fairness rises
by 0.000245. Effective airtime falls by 0.002059, while P95 delay increases by
\SI{193.3}{\micro\second}. The utility gain is therefore a collision reduction
that outweighs small airtime and delay costs; it is not a uniform improvement
of every metric.

Arm choices follow the offline conflict. Nodes choose $q_4$ for 98.9\% of
light-phase decisions. In the high phase, the fixed-TMC profile $q_{20}$ rises
to 76.9\%. Residual exploration of $q_4$ is consistent with the small high-load
penalty. This behavior also gives a protocol interpretation to the light-load
gain: queues repeatedly enter contention, and the shorter recovery cycle with
$m=10$ is favored; persistent high load returns most decisions to TMC.

We measure adaptation from the local reward trace after every explicit phase
change. Recovery is the first of three consecutive eight-decision windows that
reaches 90\% of the improvement from the immediate post-change level to the
last-quarter phase level. All 224 transitions recover without censoring.
Median $T_{\rm adapt}$ is 736 rounds and the 95th percentile is 928 rounds,
well below $T_{\rm dwell}=4096$. The experiment therefore checks, rather than
only assumes, the critical ordering in \eqref{eq:timescale} for this scenario.

\subsection{Aggregation Sensitivity and Packet Check}

A whole-run calculation pools phases before forming the nonlinear utility.
It gives $-0.000668$ ($[-0.000921,-0.000413]$), with 6/32 positive seeds.
The direction changes because the active population, delay distribution, and
fairness denominator are mixed before normalization; the corresponding
airtime, P95-delay, and fairness differences become $-0.013570$,
$+\SI{2.240}{\milli\second}$, and $-0.008170$. We retain this diagnostic because
an application that values one end-to-end mixed-population score should not
substitute the phase-conditional conclusion.

The repository also contains an ns-3.35/5G-LENA/NR-U protocol cross-check with
Wi-Fi EDCA, NR-U alignment, HARQ, and a frequency-selective channel. That older
three-seed campaign used the initial unrestricted 24-profile model, so its
mixed technology-level directions are not treated as confirmation of the
restricted controller.

\begin{table}[!t]
\caption{Five-Second ns-3 Shadowing Diagnostic (TMC, Seed 410)}
\label{tab:ns3-diagnostic}
\centering
\scriptsize
\begin{tabular}{@{}lrrrr@{}}
\toprule
Tech. & Thr. (Mb/s) & Delay (ms) & IP loss & Access overlap \\
\midrule
Wi-Fi & 8.121 & 1.793 & 0.0117 & 0.3093 \\
NR-U  & 7.537 & 38.087 & 0.0828 & 0.7726 \\
\bottomrule
\end{tabular}
\end{table}

Table~\ref{tab:ns3-diagnostic} is a measurement-path diagnostic with 3GPP
shadowing enabled, not an adaptive-policy comparison. The access quantity is
the fraction of transmission starts overlapping another active transmission;
IP loss is computed from transmitted and received packets. Their different
values confirm that decoding loss and simultaneous access must remain separate.
